\documentclass[12pt]{article}
\usepackage[utf8]{inputenc}
\usepackage{vntex} % Bộ font tiếng Việt chuẩn nhất trong LaTeX, tự động nạp T5 encoding và hỗ trợ unicode
\usepackage{amsmath}
\usepackage{tikz}
\usepackage{booktabs}
\usepackage{longtable}
\usepackage{hyperref}
\usepackage{geometry}
\usepackage{array}

% Cấu hình lề rộng rãi để bảng mô tả chi tiết chức năng không bị tràn lề
\geometry{a4paper, left=1.5cm, right=1.5cm, top=2cm, bottom=2cm}
\usetikzlibrary{shapes.geometric, arrows, positioning, shadows} 

\hypersetup{
    colorlinks=true,
    linkcolor=blue,
    filecolor=magenta,      
    urlcolor=cyan,
    pdftitle={Bao Cao Phan Quyen He Thong - StationOS},
    pdfpagemode=FullScreen,
}

% Định nghĩa phong cách vẽ TikZ
\tikzset{
    t4/.style={rectangle, rounded corners, minimum width=5cm, minimum height=1.2cm, text width=5.5cm, text centered, draw=green!60!black, fill=green!10, thick, drop shadow},
    t3/.style={rectangle, rounded corners, minimum width=5cm, minimum height=1.2cm, text width=5.5cm, text centered, draw=orange!80!black, fill=orange!10, thick, drop shadow},
    t2/.style={rectangle, rounded corners, minimum width=5cm, minimum height=1.2cm, text width=5.5cm, text centered, draw=blue!80!black, fill=blue!10, thick, drop shadow},
    t1/.style={rectangle, rounded corners, minimum width=4.5cm, minimum height=1.2cm, text width=4.5cm, text centered, draw=red!80!black, fill=red!10, thick, drop shadow},
    arrow/.style={thick, ->, >=stealth, draw=gray!80}
}

\begin{document}

\begin{titlepage}
    \centering
    \vspace*{2cm}
    {\huge\bfseries BÁO CÁO THỐNG KÊ PHÂN QUYỀN HỆ THỐNG (RBAC) - STATIONOS\par}
    \vspace{1.5cm}
    {\large\bfseries Tài liệu Đặc tả kỹ thuật \& Cấu hình phân cấp an toàn\par}
    \vspace{2cm}
    \textbf{Dự án:} Hệ thống điều hành giám sát trạm biến áp trung tâm (Master Station)\\
    \vfill
    {\large \today \par}
\end{titlepage}

\section{Giới thiệu chung}
Tài liệu này tổng hợp chi tiết các vai trò (Role), tài khoản mẫu mặc định và kiến trúc phân cấp an toàn được lập trình trong nền tảng StationOS. Mục tiêu của kiến trúc này là phân tách rạch ròi phạm vi quản lý và đảm bảo an toàn tuyệt đối cho Trạm Tổng (Master Station) trước các truy cập trái phép từ tài khoản của các Trạm Con.

\section{Sơ đồ kiến trúc phân cấp (4 Tầng Quản Trị)}
Hệ thống sử dụng mô hình phân cấp dạng cây từ cấp cao nhất xuống cấp cơ sở, cho phép lọc dữ liệu tự động theo địa bàn quản lý.

\begin{figure}[htbp]
    \centering
    \begin{tikzpicture}[node distance=1.8cm, scale=0.85, every node/.style={transform shape}]
        % Các Node đại diện cho các Tầng
        \node (T4) [t4] {\textbf{TẦNG 4: Admin Toàn Cục (multi)} \\ $\bullet$ Quản trị tối cao (Trạm Tổng) \\ $\bullet$ Quản lý không giới hạn Tỉnh/Trạm};
        
        \node (T3_LA) [t3, below left=1.5cm and 1.5cm of T4] {\textbf{TẦNG 3: Admin Tỉnh (provinceadmin)} \\ $\bullet$ Quản lý Tỉnh Long An \\ $\bullet$ Quản lý các Tổ/Trạm thuộc Long An};
        
        \node (T3_TN) [t3, below right=1.5cm and 1.5cm of T4] {\textbf{TẦNG 3: Admin Tỉnh (provinceadmin)} \\ $\bullet$ Quản lý Tỉnh Tây Ninh \\ $\bullet$ Quản lý các Tổ/Trạm thuộc Tây Ninh};
        
        \node (T2_LA) [t2, below=of T3_LA] {\textbf{TẦNG 2: Tổ trưởng (teamleader)} \\ $\bullet$ Phụ trách Tổ 1 (Long An) \\ $\bullet$ Giám sát các Trạm A, B, C};
        
        \node (T2_TN) [t2, below=of T3_TN] {\textbf{TẦNG 2: Tổ trưởng (teamleader)} \\ $\bullet$ Phụ trách Tổ 2 (Tây Ninh) \\ $\bullet$ Giám sát các Trạm D, E};
        
        \node (T1_LA) [t1, below=of T2_LA] {\textbf{TẦNG 1: Admin Trạm (stationadmin)} \\ $\bullet$ Quản lý cục bộ Trạm Long An};
        
        \node (T1_TN) [t1, below=of T2_TN] {\textbf{TẦNG 1: Admin Trạm (stationadmin)} \\ $\bullet$ Quản lý cục bộ Trạm Tây Ninh};

        % Đường nối kết nối mối quan hệ
        \draw [arrow] (T4) -- (T3_LA);
        \draw [arrow] (T4) -- (T3_TN);
        \draw [arrow] (T3_LA) -- (T2_LA);
        \draw [arrow] (T3_TN) -- (T2_TN);
        \draw [arrow] (T2_LA) -- (T1_LA);
        \draw [arrow] (T2_TN) -- (T1_TN);
    \end{tikzpicture}
    \caption{Sơ đồ cây phân cấp quản trị 4 tầng trong hệ thống StationOS.}
    \label{fig:rbac_tree}
\end{figure}

\newpage

\section{Bảng Phân Quyền Chi Tiết các Tài Khoản Mặc Định}
Hệ thống hỗ trợ 4 tài khoản quản trị mặc định chính tương ứng với các tầng phân cấp:

\begin{longtable}{>{\centering\arraybackslash}p{0.6cm} p{2.2cm} p{2.0cm} p{2.2cm} p{3.2cm} p{7.3cm}}
\caption{Danh sách tài khoản quản trị mặc định và chức năng chi tiết.} \label{tab:rbac} \\
\toprule
\textbf{STT} & \textbf{Cấp Quản Lý} & \textbf{Tài Khoản (Username)} & \textbf{Mật Khẩu Mặc Định} & \textbf{Phạm Vi (Scope)} & \textbf{Chi Tiết Chức Năng Cho Phép} \\
\midrule
\endfirsthead
\multicolumn{6}{c}%
{{\bfseries Bảng \thetable\ tiếp theo từ trang trước}} \\
\toprule
\textbf{STT} & \textbf{Cấp Quản Lý} & \textbf{Tài Khoản} & \textbf{Mật Khẩu} & \textbf{Phạm Vi (Scope)} & \textbf{Chi Tiết Chức Năng Cho Phép} \\
\midrule
\endhead
\bottomrule
\endfoot
\bottomrule
\endlastfoot

1 & Admin Toàn Cục & \texttt{multi} & \texttt{Demo@2024} & Toàn bộ hệ thống (Không giới hạn Tỉnh/Trạm). & 
\begin{itemize}
    \item Thêm, sửa, xóa các Tỉnh/Đơn vị thành phố trong hệ thống.
    \item Thêm mới, cấu hình kết nối, gỡ liên kết các Trạm Con vào Tỉnh.
    \item Tạo mới và phân quyền cho tài khoản Admin cấp dưới (\texttt{provinceadmin}, \texttt{teamleader}, \texttt{stationadmin}).
    \item Quản lý, kích hoạt, xem giới hạn Giftcode bản quyền toàn cục.
    \item Giám sát camera, biểu đồ đo đạc, lịch sử cảnh báo và báo cáo phân tích trên toàn bộ hệ thống.
    \item Điều khiển thiết bị từ xa tại bất kỳ trạm con nào.
\end{itemize} \\
\midrule
2 & Admin Tỉnh & \texttt{provinceadmin} & \texttt{Province@123} & Chỉ trong Tỉnh được gán (Ví dụ: Tây Ninh, Long An...). & 
\begin{itemize}
    \item Quản lý danh sách các Trạm Con và các Tổ thao tác lưu động trực thuộc tỉnh được phân công.
    \item Tạo mới và quản lý tài khoản nhân sự cấp dưới (Tổ trưởng, Admin Trạm, Nhân viên vận hành) thuộc tỉnh.
    \item Kích hoạt, gia hạn Giftcode bản quyền cho các Trạm Con trong tỉnh.
    \item Xem trực tuyến camera, bản đồ số, trạng thái thiết bị và báo cáo tổng hợp thuộc địa bàn quản lý.
\end{itemize} \\
\midrule
3 & Tổ trưởng Tổ thao tác & \texttt{teamleader} & \texttt{TeamLeader@123} & Các trạm do Tổ thao tác phụ trách. & 
\begin{itemize}
    \item Giám sát thời gian thực các trạm biến áp được phân công phụ trách.
    \item Xem trực tiếp camera, nhận cảnh báo cháy, nhiệt độ cao, xâm nhập khẩn cấp tại hiện trường để đi xử lý sự cố.
    \item Điều khiển bật/tắt thiết bị ngoại vi (còi, đèn cảnh báo), điều khiển quay quét camera PTZ khi làm nhiệm vụ.
    \item Xem nhật ký sự cố và xuất báo cáo vận hành của các trạm thuộc tổ quản lý.
\end{itemize} \\
\midrule
4 & Admin Trạm & \texttt{stationadmin} & \texttt{Station@123} & Chỉ trong Trạm con được gán. & 
\begin{itemize}
    \item Thiết lập sơ đồ kết nối camera, cảm biến nhiệt độ, cảm biến PD cục bộ tại trạm con.
    \item Cấu hình các luật cảnh báo tự động tại trạm (ngưỡng nhiệt độ an toàn, vùng quét phát hiện xâm nhập).
    \item Quản lý và phân công lịch trực cho nhân viên vận hành thuộc trạm.
    \item Xem camera, lịch sử ghi hình, nhật ký cảnh báo và báo cáo hiệu suất thiết bị của trạm con.
\end{itemize} \\
\end{longtable}

\section{Các Quy Tắc Bảo Mật Đặc Biệt Trong Mã Nguồn}
Nhằm đảm bảo an ninh tuyệt đối cho hệ thống, các quy tắc bảo mật sau đã được áp dụng cứng trong backend:

\begin{itemize}
    \item \textbf{Tách biệt an toàn Trạm Tổng \& Trạm Con}: Tài khoản \texttt{admin} (mật khẩu \texttt{Admin@123}) chỉ tồn tại ở các Trạm Con nhằm phục vụ giao thức API đồng bộ dữ liệu. Trên cơ sở dữ liệu của Trạm Tổng (\texttt{StationOS\_Central}), tài khoản này đã được \textbf{xóa hoàn toàn} nhằm chặn đứng nguy cơ dùng tài khoản này đăng nhập trái phép vào Trạm Tổng. (Đã sửa ký tự gạch dưới để không lỗi biên dịch).
    \item \textbf{Bảo vệ tài khoản mặc định}: Các tài khoản admin hệ thống mặc định (\texttt{multi}, \texttt{provinceadmin}, \texttt{stationadmin}, \texttt{teamleader}) sẽ tự động bị ẩn các chức năng Sửa, Đổi mật khẩu, Vô hiệu hóa trên giao diện để tránh trường hợp quản trị viên vô tình tự khóa tài khoản của mình.
    \item \textbf{Đồng bộ hóa thời gian thực (Realtime Update)}: Mọi sự thay đổi về gán quyền hoặc sửa đổi Tổ thao tác sẽ được truyền phát ngay lập tức tới các Client đang kết nối qua SignalR, giúp cập nhật giao diện tự động mà không cần người dùng tải lại trang.
\end{itemize}

\end{document}
